\documentclass[11pt,a4paper]{article}
\usepackage[margin=25mm]{geometry}
\usepackage{amsmath,amssymb,amsthm,booktabs}
\newtheorem{proposition}{Proposition}
\newtheorem{lemma}[proposition]{Lemma}
\title{Rotation conventions and closed flag walks}
\author{Human Compression project}\date{}
\begin{document}\maketitle
The survey of Schulte--Weiss uses consecutive rotations, whereas our
chamber arguments use rotations based at one fixed type. We give the
change of generators and recover the survey's period formulas for holes
and zigzags. The formulas are established results recorded in the survey;
no new classification or answer to its Problem~7 is asserted.

\section{Two choices of rotations}
Let $n\geq3$ and let a group $G$ have generators $\sigma_1,\ldots,\sigma_{n-1}$ satisfying
\[
 (\sigma_i\sigma_{i+1}\cdots\sigma_j)^2=1
 \qquad(1\leq i<j\leq n-1).
\]
Put $\alpha_0=1$ and $\alpha_i=\sigma_1\cdots\sigma_i$. Then
$\sigma_i=\alpha_{i-1}^{-1}\alpha_i$, and $\alpha_i^2=1$ for $i\geq2$.
These are identities in the marked group, independent of an assertion
that its coset geometry is a polytope.

\begin{proposition}
An automorphism $\theta$ simultaneously inverts every $\alpha_i$ if and
only if
\[
 \theta(\sigma_1)=\sigma_1^{-1},\qquad
 \theta(\sigma_2)=\sigma_1^2\sigma_2,\qquad
 \theta(\sigma_i)=\sigma_i\quad(i\geq3).
\]
If such an automorphism exists, its square is the identity.
\end{proposition}
\begin{proof}
Suppose first that $\theta(\alpha_i)=\alpha_i^{-1}$. The assertion for
$\sigma_1=\alpha_1$ follows at once. Since $\alpha_2$ is an involution,
\[
 \theta(\sigma_2)=\theta(\alpha_1^{-1}\alpha_2)
 =\alpha_1\alpha_2=\sigma_1^2\sigma_2.
\]
For $i\geq3$ both $\alpha_{i-1}$ and $\alpha_i$ are involutions, so they
are fixed by $\theta$, and their product $\sigma_i$ is fixed.
Conversely the displayed formulas give
$\theta(\alpha_1)=\alpha_1^{-1}$ and
$\theta(\alpha_2)=\sigma_1^{-1}\sigma_1^2\sigma_2=\alpha_2$.
Multiplying by the fixed later generators gives
$\theta(\alpha_i)=\alpha_i=\alpha_i^{-1}$ for every $i\geq2$.
Applying $\theta$ twice fixes the generating set $(\alpha_i)$.
\end{proof}
Thus the two reflection tests agree whenever the appropriate
string-group reconstruction theorem applies. We retain that reconstruction
theorem, including its intersection hypotheses, as an attributed input.
A presentation with the displayed relations alone is not a reconstruction
proof. Also $\theta$ may be the identity; in that case a geometric
reflection interchanging two chamber orbits is still a separate
involution, as in the accepted residue theorem.

\section{Automorphisms and adjacency operations}
For a polyhedron let $r_i$ send a flag to its unique $i$-adjacent flag.
The symbols $r_i$ in this section denote permutations of flags. They
need not be automorphisms of the polyhedron. They satisfy $r_i^2=1$
and $r_0r_2=r_2r_0$. A word is applied in its written order; write
$\Phi^w$ for its result on a flag $\Phi$. Automorphisms commute with
all these operations because adjacency is intrinsic.

\begin{lemma}\label{period}
Suppose the automorphism group acts freely and transitively on a set
of flags $\mathcal O$, and a word $w$ preserves $\mathcal O$.
Let $g_w$ be the unique automorphism with $g_w\Phi=\Phi^w$.
Then $\Phi^{w^k}=g_w^k\Phi$. The least positive return time of this
flag walk is the order of $g_w$, with both interpreted as infinite
if there is no return.
\end{lemma}
\begin{proof}
An automorphism fixing one flag fixes its adjacent flags and hence all
flags by connectedness; thus the freeness hypothesis holds in every
application here. Commutation gives
$(g_w^k\Phi)^w=g_w^k(\Phi^w)=g_w^{k+1}\Phi$, proving the identity by
induction. The final assertion is freeness.
\end{proof}
In a regular polyhedron $\mathcal O$ is the full flag set. In a chiral
polyhedron there are two such orbits and every adjacency interchanges
them. A word of even length preserves each orbit, and a word of odd
length interchanges them. Accordingly the lemma applies to even words
in the chiral case.

\section{Holes and zigzags}
A flag records a vertex, an incident edge and one incident face.
Crossing the edge is the operation $r_0$. At the new vertex, the
operation $r_1$ takes the next edge in the selected face, and $r_2$
changes the selected face while keeping the vertex and edge. Repeating
these two local operations goes around the vertex figure.

Fix $1\leq j\leq\lfloor q/2\rfloor$ for a polyhedron of type $\{p,q\}$.
The following words each advance by one edge of the indicated walk:
\[
 h_j=r_0r_1(r_2r_1)^{j-1},\qquad
 z_j=r_0(r_1r_2)^j.
\]
Indeed after $r_0$, the first $r_1$ selects the first outgoing edge;
each additional $r_2r_1$ selects the next one. This gives the $j$th
outgoing edge. In $h_j$ the selected face remains on the same side of
the walk, and the successive turns have the same sense. The final
$r_2$ in $z_j=h_jr_2$ chooses the opposite incident face, so the next
turn has the opposite sense. These are respectively the $j$-hole and
$j$-zigzag rules. Their lengths are their complete traversal periods,
including the selected side of the traversal, rather than the first
return to a previously visited vertex.

For a regular polyhedron, let $\rho_i$ be the distinguished
automorphism with $\rho_i\Phi=\Phi^{r_i}$. Commutation with adjacency
shows inductively that a word in the $\rho_i$ has the same effect on
$\Phi$ as the written word in the $r_i$. Lemma~\ref{period} therefore
gives
\[
 \ell(h_j)=\operatorname{ord}\bigl(\rho_0\rho_1
                (\rho_2\rho_1)^{j-1}\bigr),\qquad
 \ell(z_j)=\operatorname{ord}\bigl(\rho_0(\rho_1\rho_2)^j\bigr).
\]

For a chiral polyhedron choose the distinguished rotations by
$\sigma_1\Phi=\Phi^{r_0r_1}$ and
$\sigma_2\Phi=\Phi^{r_1r_2}$. These are consecutive rotations with the
chosen base orientation. Write $s_1=r_0r_1$, $s_2=r_1r_2$ as formal
words in involutions. Direct cancellation gives
\[
 h_j=s_1s_2^{1-j},\qquad
 r_0s_2^jr_0=s_1s_2^{-j}s_1^{-1},\qquad
 z_j^2=s_1s_2^{-j}s_1^{-1}s_2^j.
\]
For $j=1$, direct cancellation gives
\[
 (r_0r_1)(r_2r_1)(r_1r_0)=r_0r_1r_2r_0=r_0s_2r_0.
\]
Taking the $j$th power proves the middle identity.
The first and last words are even. Applying the lemma on the base
flag orbit proves
\begin{align*}
 \ell(h_j)&=\operatorname{ord}(\sigma_1\sigma_2^{1-j}),\\
 \ell(z_j)&=2\operatorname{ord}
       (\sigma_1\sigma_2^{-j}\sigma_1^{-1}\sigma_2^j).
\end{align*}
The factor two is necessary: one zigzag step exchanges the two flag
orbits, so a return cannot have odd length. The alternative word
$\sigma_2^j\sigma_1\sigma_2^{-j}\sigma_1^{-1}$ is conjugate to the
word displayed above, by $\sigma_2^j$, and has the same order.

These computations identify the exact flag periods. They neither say
that the period data determine the polyhedron nor prove finiteness of
a universal quotient with prescribed periods. Those are the separate
questions posed in the survey.

\begin{thebibliography}{9}\small
\bibitem{survey} E. Schulte and A. Ivi\'c Weiss,
\emph{Problems on polytopes, their groups, and realizations},
Period. Math. Hungar. 53 (2006), 231--255; fixed arXiv math/0608397v1,
Sections 2--3. We use the survey's definitions and its explicit attribution
of the reconstruction theorem to their earlier work on chiral polytopes.
\end{thebibliography}
\end{document}
